\chapter{Lipschitz maps acting on compact normal currents}\label{sec:normal-lipschitz-pushforward}
The homotopy estimate gives an analytic definition of a Lipschitz pushforward before any integer-rectifiability theorem is used. We develop the compact-support case of \cite{Sim18}, Chapter 6, Lemmas 2.35--2.36, printed pp.~166--167. All norms on forms are comass, and $\mathcal F_{\mathbb R}$ is the real-current flat norm from Definition \ref{def:real-current-flat-norm}. Preservation of integer rectifiability and identification with a geometric tangential formula are developed separately in Section \ref{sec:geometric-integer-pushforward}; the results below do not assume them.

\begin{lemma}[Vector mollification preserves the operator Lipschitz bound]\label{lem:vector-map-lipschitz-mollification}
Let $f:\mathbb R^N\to\mathbb R^Q$ be $L$-Lipschitz. For a nonnegative smooth mollifier $\rho$ supported in the unit ball with integral one, set $f_\varepsilon=f*\rho_\varepsilon$. Then $f_\varepsilon$ is smooth,
\[
 \sup_x|f_\varepsilon(x)-f(x)|
 \leq L\varepsilon\int|z|\rho(z)\,dz,
 \qquad \operatorname{Lip}f_\varepsilon\leq L,
 \qquad \|Df_\varepsilon\|_{\rm op}\leq L.
\]
There is no factor depending on the target dimension in this bound.
\end{lemma}
\begin{proof}
Write the convolution as the vector integral $f_\varepsilon(x)=\int f(x-\varepsilon z)\rho(z)\,dz$. Its difference from $f(x)$ is bounded by the integral of $L\varepsilon|z|$. For two points $x,y$, estimate the norm of the difference of the two integrals by $L|x-y|\int\rho=L|x-y|$. Smoothness follows from local integrability and differentiation of the compactly supported kernel. Difference quotients of the smooth map give the operator bound. The norm is estimated on the vector integral itself, so a coordinatewise extension estimate with a square-root dimension factor is unnecessary. This is the vector version of Proposition \ref{prop:lipschitz-mollification}.
\end{proof}

\begin{theorem}[Analytic pushforward of a compact normal current]\label{thm:compact-normal-lipschitz-pushforward}
Let $T$ be a compactly supported normal $m$-current in $\mathbb R^N$, $m\geq0$, and let $f:\mathbb R^N\to\mathbb R^Q$ be $L$-Lipschitz. For every smooth compactly supported $m$-form $\omega$, the limit
\[
 f_\#T(\omega)=\lim_{\varepsilon\downarrow0}(f_\varepsilon)_\#T(\omega)
\]
exists and defines a compactly supported normal current. It satisfies
\[
 \operatorname{spt}(f_\#T)\subset f(\operatorname{spt}T),
 \quad \mathbf M(f_\#T)\leq L^m\mathbf M(T),
 \quad \partial(f_\#T)=f_\#(\partial T).
\]
For $m\geq1$ its boundary mass is at most $L^{m-1}\mathbf M(\partial T)$. Use $L^0=1$, including $L=0$; a zero-current has zero boundary. This is the compact-support form of \cite{Sim18}, Lemma 2.36.
\end{theorem}
\begin{proof}
In dimension zero, Theorem \ref{thm:current-mass-representation} identifies $T$ with a finite signed measure. Smooth pushforward tests it against $\omega\circ f_\varepsilon$, which converges uniformly on its compact support to $\omega\circ f$. This gives the usual measure pushforward, its mass bound and support assertion. Its boundary is zero.

For $m\geq1$, put $K=\operatorname{spt}T$. Proposition \ref{prop:current-homotopy-estimate} applied to $f_\varepsilon,f_\delta$ gives
\[
 \mathcal F_{\mathbb R}((f_\varepsilon)_\#T-(f_\delta)_\#T)
 \leq \sup_K|f_\varepsilon-f_\delta|
       \bigl(L^m\mathbf M(T)+L^{m-1}\mathbf M(\partial T)\bigr).
\]
The right side tends to zero by Lemma \ref{lem:vector-map-lipschitz-mollification}. Proposition \ref{prop:real-flat-control} makes evaluation on each fixed test form Cauchy. Define the limiting linear functional by these evaluations. Proposition \ref{prop:current-pushforward-mass} gives $|f_\#T(\omega)|\leq L^m\mathbf M(T)\|\omega\|_{*,\infty}$, so the functional is a current with that mass bound.

If a test form is supported off the compact set $f(K)$, its support has positive distance from $f(K)$. Uniform convergence on $K$ makes its smooth pullback vanish near $K$ for sufficiently small parameters. Hence the limit has support contained in $f(K)$. For a test $(m-1)$-form $\eta$, Proposition \ref{prop:current-pushforward-identities} gives
\[
 \partial f_\#T(\eta)
 =\lim_\varepsilon (f_\varepsilon)_\#T(d\eta)
 =\lim_\varepsilon (f_\varepsilon)_\#(\partial T)(\eta)
 =f_\#(\partial T)(\eta).
\]
The construction for $\partial T$ is already valid by induction on dimension; it is a compact normal current because $\partial\partial T=0$. Its mass bound proves the boundary estimate and normality. We used limits on tests and mass bounds, not an unproved completeness assertion for the real flat norm or an integral-current compactness theorem.
\end{proof}

\begin{proposition}[Independence of approximation and agreement on the support]\label{prop:normal-pushforward-approximation-independence}
Let $T,f$ be as in Theorem \ref{thm:compact-normal-lipschitz-pushforward}. Suppose smooth maps $g_j$ converge to $f$ uniformly on $K=\operatorname{spt}T$ and have a common derivative bound on one fixed neighborhood of $K$. Then $(g_j)_\#T\rightharpoonup f_\#T$. In particular the definition is independent of the mollifier. If two globally Lipschitz maps agree on $K$, their pushforwards of $T$ agree. This is the compact-support content of \cite{Sim18}, Lemma 2.35 and the argument for Lemma 2.36.
\end{proposition}
\begin{proof}
For $m\geq1$, compare $g_j$ with $f_{1/j}$. Their derivative bounds near $K$ are bounded by one common number $L_0$, and their uniform difference on $K$ tends to zero. The homotopy estimate and flat control of test forms give the same weak limit. For dimension zero, uniform convergence of the scalar tests composed with these maps gives the conclusion directly through the measure representation. Applying this argument to two mollifier choices proves independence.

If $f=g$ on $K$, compare their mollifications. On $K$, their difference tends uniformly to zero, and both derivative bounds stay uniformly bounded. The same estimate shows that their pushforwards agree on every test form. This uses normality through the boundary mass in the homotopy estimate; finite mass alone does not supply it.
\end{proof}

\begin{proposition}[Flat stability under a change of Lipschitz map]\label{prop:normal-lipschitz-map-flat-stability}
Let $m\geq1$ and $T$ be compactly supported and normal. If $f,g:\mathbb R^N\to\mathbb R^Q$ are $L$-Lipschitz and $\delta=\sup_{\operatorname{spt}T}|f-g|$, then
\[
 \mathcal F_{\mathbb R}(g_\#T-f_\#T)
 \leq\delta\bigl(L^m\mathbf M(T)+L^{m-1}\mathbf M(\partial T)\bigr).
\]
In particular the smooth approximations in Theorem \ref{thm:compact-normal-lipschitz-pushforward} converge in the real flat norm, as do pushforwards by uniformly convergent maps with a common global Lipschitz bound.
\end{proposition}
\begin{proof}
For smooth mollifications use the affine homotopy. Its filling $S_\varepsilon$ and residual $R_\varepsilon$ satisfy
\[
 (g_\varepsilon)_\#T-(f_\varepsilon)_\#T
 =R_\varepsilon+\partial S_\varepsilon,
\]
with $\mathbf M(S_\varepsilon)\leq\delta_\varepsilon L^m\mathbf M(T)$ and $\mathbf M(R_\varepsilon)\leq\delta_\varepsilon L^{m-1}\mathbf M(\partial T)$, where $\delta_\varepsilon=\sup_K|f_\varepsilon-g_\varepsilon|\to\delta$. All residuals and fillings have support in one compact ball: the images of $K$ under both map sequences are uniformly bounded, and the affine homotopies lie in their convex hull.

Apply Proposition \ref{prop:current-weak-compactness} simultaneously to these two sequences. Extract weak limits $R,S$; their supports remain in the same compact ball. Taking limits on test forms and their exterior derivatives gives $g_\#T-f_\#T=R+\partial S$. Mass is lower semicontinuous under weak convergence, directly by taking the supremum of limiting evaluations on forms of comass at most one. Thus the limiting mass bounds are at most the stated bounds with $\delta$. Insert $R,S$ into the definition of $\mathcal F_{\mathbb R}$.

For a smooth map its analytic pushforward agrees with the original smooth one by Proposition \ref{prop:normal-pushforward-approximation-independence}, taking the constant smooth approximating sequence. The bound applied to $f_\varepsilon,f$ therefore proves the last assertion. No lower semicontinuity theorem for the flat norm itself is needed; explicit weak limits of the filling and residual give admissible competitors.
\end{proof}

\begin{lemma}[A smooth globally Lipschitz map bounds the real flat norm]\label{lem:smooth-map-real-flat-bound}
Let $h:\mathbb R^Q\to\mathbb R^P$ be smooth and $L$-Lipschitz, and let $X$ be a compactly supported $m$-current of finite real flat norm. Then
\[
 \mathcal F_{\mathbb R}(h_\#X)
 \leq\max\{L^m,L^{m+1}\}\,\mathcal F_{\mathbb R}(X).
\]
All maps are proper on the compact supports of the currents used in the infimum.
\end{lemma}
\begin{proof}
For any admissible decomposition $X=R+\partial S$, smooth pushforward and boundary commutation give $h_\#X=h_\#R+\partial h_\#S$. Their masses are at most $L^m\mathbf M(R)$ and $L^{m+1}\mathbf M(S)$ by Proposition \ref{prop:current-pushforward-mass}. Their supports are compact, being contained in continuous images of compact supports. They are therefore admissible competitors for the target flat norm. Add the bounds and take the infimum. This lemma is stated for smooth maps, so it does not presuppose a Lipschitz pushforward for arbitrary finite-mass competitors that need not be normal.
\end{proof}

\begin{proposition}[Composition for compact normal Lipschitz pushforwards]\label{prop:normal-lipschitz-pushforward-composition}
For a compactly supported normal $m$-current $T$ and globally Lipschitz maps $f:\mathbb R^N\to\mathbb R^Q$ and $g:\mathbb R^Q\to\mathbb R^P$,
\[
 (g\circ f)_\#T=g_\#(f_\#T).
\]
The intermediate pushforward is compactly supported and normal, so the right side is defined by Theorem \ref{thm:compact-normal-lipschitz-pushforward}.
\end{proposition}
\begin{proof}
For $m=0$, this is composition of pushforwards of finite signed measures. Let $m\geq1$ and mollify both maps to $f_j,g_j$, retaining their global Lipschitz bounds. The maps $g_j\circ f_j$ converge uniformly to $g\circ f$ on $\operatorname{spt}T$: split the error into the uniform approximation of $g_j$ and the Lipschitz bound for $g$ times the error of $f_j$. Their derivatives have one common bound. Proposition \ref{prop:normal-pushforward-approximation-independence} identifies the weak limit of $(g_j\circ f_j)_\#T$ with $(g\circ f)_\#T$.

By Proposition \ref{prop:normal-lipschitz-map-flat-stability}, $(f_j)_\#T\to f_\#T$ in the real flat norm. Lemma \ref{lem:smooth-map-real-flat-bound} therefore gives
\[
 \mathcal F_{\mathbb R}\bigl((g_j)_\#((f_j)_\#T)-(g_j)_\#(f_\#T)\bigr)\to0,
\]
since the Lipschitz bounds of $g_j$ are uniform. Apply Proposition \ref{prop:normal-lipschitz-map-flat-stability} again, now to the compact normal current $f_\#T$, to obtain $(g_j)_\#(f_\#T)\to g_\#(f_\#T)$ in flat norm. Proposition \ref{prop:current-pushforward-identities} identifies the first term with $(g_j\circ f_j)_\#T$. Flat control of test forms identifies the two weak limits and proves the formula.
\end{proof}

\begin{example}[Finite mass without normality does not determine pushforward from support values]\label{ex:transverse-current-support-values}
Define a compactly supported one-current in $\mathbb R^2$ by
\[
 T(a\,dx+b\,dy)=\int_0^1 b(x,0)\,dx.
\]
It has mass one and support $[0,1]\times\{0\}$, but is not normal. The smooth maps $f(x,y)=0$ and $g(x,y)=y$ agree on its support, while $f_\#T=0$ and
\[
 g_\#T(\phi(t)\,dt)=\phi(0).
\]
Thus normality is essential in the support-agreement assertion above.
\end{example}
\begin{proof}
The polar vector is the unit vertical vector $e_y$ and its measure is length on the horizontal segment, giving mass one. Testing nonnegative $dy$ coefficients localized near any point of the segment gives the asserted support. For a compactly supported smooth scalar test $\eta$,
\[
 \partial T(\eta)=\int_0^1\partial_y\eta(x,0)\,dx.
\]
Choose a smooth compactly supported function $\chi(x)$ equal to one near $[0,1]$, and a smooth cutoff $\zeta(y)$ equal to one near zero, both between zero and one. The tests $\eta_j(x,y)=\chi(x)\zeta(y)\sin(jy)$ have supremum norm at most one, while $\partial T(\eta_j)=j$. Thus the boundary does not have finite mass. Pullback by the constant map annihilates every one-form; pullback by $g$ gives $\phi(y)\,dy$, proving its evaluation. The polar direction is transverse to the support and is not the approximate tangent orientation of a rectifiable one-current.
\end{proof}

\begin{example}[A nonproper parametrization can make the proposed pushforward infinite]\label{ex:nonproper-periodic-current-map}
Let $T=\llbracket\mathbb R\rrbracket$ be the integration current on the oriented line and let $f(t)=(\cos t,\sin t)$. The current is locally integral and the map is globally Lipschitz, but $f$ is not proper on the support. Its formal pullback action on one compactly supported target form has infinite integral, so that action does not define a current.
\end{example}
\begin{proof}
Choose a smooth compactly supported cutoff $\chi$ on $\mathbb R^2$ which equals one near the unit circle, and put $\omega=\chi(x,y)(x\,dy-y\,dx)$. Along $f$, its pullback is $dt$, because $\cos^2t+\sin^2t=1$. The integral on $[-R,R]$ equals $2R$ and tends to infinity. Thus the whole-line integral is not a finite test evaluation. The inverse image of the compact unit circle is all of $\mathbb R$, verifying failure of properness. This does not contradict the compact-support construction: $T$ here has noncompact support and infinite total mass.
\end{proof}
